% This survey highlights the transformative potential of large language models and agentic workflows for automating high-performance kernel generation, signaling a paradigm shift from static code synthesis toward dynamic, feedback-driven optimization. By synthesizing recent advances in supervised fine-tuning, reinforcement learning, and multi-agent orchestration, together with progress in kernel-centric dataset curation and benchmark development, we present a unified perspective on this rapidly evolving research area. Looking ahead, future work must move beyond rigid workflow-based designs toward deeper agentic reasoning and self-evolving optimization capabilities, while simultaneously addressing ecosystem bias to support diverse hardware platforms and broadly generalizable kernel-generation methodologies. Ultimately, addressing these challenges is critical not only for reducing the burden of manual kernel engineering and freeing human experts from low-level optimization complexity, but also for systematically improving software efficiency and unlocking productivity gains in the era of rapidly scaling AI infrastructure.


This survey highlights the transformative potential of large language models and agentic workflows for automating high-performance kernel generation, synthesizing recent advances in supervised fine-tuning, reinforcement learning, and multi-agent orchestration, together with progress in kernel-centric dataset and benchmark development. Looking ahead, future work should move beyond rigid workflows toward self-evolving agentic reasoning with strong hardware generalization. Such a shift is essential not only to alleviate the burden of manual kernel engineering, but also to unlock substantial productivity gains in the face of rapidly scaling AI infrastructure


